%==============================================================================
% Sjabloon onderzoeksvoorstel bachproef
%==============================================================================
% Gebaseerd op document class `hogent-article'
% zie <https://github.com/HoGentTIN/latex-hogent-article>

% Voor een voorstel in het Engels: voeg de documentclass-optie [english] toe.
% Let op: kan enkel na toestemming van de bachelorproefcoördinator!
\documentclass{hogent-article}

% Invoegen bibliografiebestand
\addbibresource{voorstel.bib}

% Informatie over de opleiding, het vak en soort opdracht
\studyprogramme{Professionele bachelor toegepaste informatica}
\course{Bachelorproef}
\assignmenttype{Onderzoeksvoorstel}
% Voor een voorstel in het Engels, haal de volgende 3 regels uit commentaar
% \studyprogramme{Bachelor of applied information technology}
% \course{Bachelor thesis}
% \assignmenttype{Research proposal}

\academicyear{2026-2027} % TODO: pas het academiejaar aan

% TODO: Werktitel
\title{Optimalisatie en GDPR-conforme Cloud-centralisatie van Zorglab Simulatiedata op Microsoft Azure}

% TODO: Studentnaam en emailadres invullen
\author{Gaétan Dessomviele}
\email{gaetan.dessomviele@student.hogent.be}

\projectrepo{https://github.com/ProProgrammerHassan/BachelorproefGaetanDessomviele}

% TODO: Medestudent
% Gaat het om een bachelorproef in samenwerking met een student in een andere
% opleiding? Geef dan de naam en emailadres hier
% \author{Yasmine Alaoui (naam opleiding)}
% \email{yasmine.alaoui@student.hogent.be}

% TODO: Geef de co-promotor op
%\supervisor[Co-promotor]{S. Beekman (Synalco, \href{mailto:sigrid.beekman@synalco.be}{sigrid.beekman@synalco.be})}

% Binnen welke specialisatierichting uit 3TI situeert dit onderzoek zich?
% Kies uit deze lijst:
%
% - Mobile \& Enterprise development
% - AI \& Data Engineering
% - Functional \& Business Analysis
% - System \& Network Administrator
% - Mainframe Expert
% - Als het onderzoek niet past binnen een van deze domeinen specifieer je deze
%   zelf
%
\specialisation{AI \& Data Engineering}
\keywords{GDPR, Cloud-centralisatie, Zorglab, Simulatiedata, Microsoft Azure} % TODO

\begin{document}

\begin{abstract}
  Efficiënte opslag en beheer van simulatiedata vormen een cruciale uitdaging voor het Zorglab van HOGENT, waar grote hoeveelheden videodata binnenkomen via Noldus VISO, Axis-camera's en Microsoft Teams. Het ontbreken van een gestructureerde cloud-omgeving resulteert enerzijds in onnodige opslagkosten en schaalbaarheidsproblemen, en anderzijds in risico's op het vlak van privacy en compliance. Dit onderzoek beantwoordt de vraag: hoe kan een kostenefficiënte en GDPR-conforme cloudarchitectuur op Microsoft Azure worden ontworpen en als Proof of Concept worden uitgewerkt voor de opslag van Zorglab-simulatiedata, rekening houdend met dataminimalisatie en regelgebaseerde filtering? De methodologie bevat vier iteratieve fasen: een grondige analyse van het privacykader (GDPR, EU-datalocatie en DPO-richtlijnen), het in kaart brengen van de huidige datastromen, het bepalen van regelgebaseerde filterregels om data efficiënt over Azure-opslaglagen te verdelen en tot slot het ontwikkelen van een functioneel Proof of Concept. Het ontwikkelde model kan het Zorglab ondersteunen bij het optimaliseren van de data-retentie, wat leidt tot strikte naleving van de privacywetgeving, gereduceerde opslagkosten en een schaalbare basis voor toekomstig onderwijskundig en wetenschappelijk onderzoek.
\end{abstract}

\tableofcontents

% De hoofdtekst van het voorstel zit in een apart bestand, zodat het makkelijk
% kan opgenomen worden in de bijlagen van de bachelorproef zelf.
%---------- Inleiding ---------------------------------------------------------

% TODO: Is dit voorstel gebaseerd op een paper van Research Methods die je
% vorig jaar hebt ingediend? Heb je daarbij eventueel samengewerkt met een
% andere student?
% Zo ja, haal dan de tekst hieronder uit commentaar en pas aan.


\section{Inleiding}%
\label{sec:inleiding}

Het Zorglab van HOGENT is een geavanceerde simulatieomgeving waarin studenten gezondheidszorg en onderzoekers praktijksituaties oefenen en analyseren. Om leerprocessen te ondersteunen en onderzoeksdata te verzamelen, wordt er uitgebreid gebruikgemaakt van video- recordings. De huidige technische opzet steunt op apparatuur en software zoals Noldus VISO, gecombineerd met Axis-camera’s en Microsoft Teams.

Deze werkwijze brengt echter belangrijke uitdagingen met zich mee. Ten eerste zorgen de hoogwaardige videobestanden voor snelgroeiende datavolumes, wat druk legt op de lokale opslaginfrastructuur en de beheerskosten verhoogt. Ten tweede – en dat is het voornaamste knelpunt – zijn videobeelden van studenten en simulanten aan te merken als gevoelige persoonsgegevens. Onder de Algemene Verordening Gegevensbescherming (AVG/GDPR) geldt een streng juridisch kader rond opslaglocaties, toegangsbeheer, bewaartermijnen en het expliciete onderscheid tussen onderwijsdata (tijdelijk) en onderzoeksdata (langduriger, met specifieke toestemming).

Bovendien mag niet zomaar alle data onbeperkt worden opgeslagen: het principe van dataminimalisatie vereist dat beelden kritisch worden gefilterd en dat overbodige of niet-relevante opnames niet doorstromen naar langdurige opslag.

\subsection{Doelstelling \& Relevantie}%
\label{subsec:doelstelling}

Het doel van deze bachelorproef is het ontwerpen en demonstreren van een veilige, geautomatiseerde en kostenefficiënte cloudarchitectuur op Microsoft Azure. Door gebruik te maken van Azure Blob Storage Tiers en geautomatiseerde opslag- en opschoonregels wordt een oplossing geboden die voldoet aan de strenge privacy-eisen en de opslagkosten formuleert naar daadwerkelijk gebruik.

\subsection{Onderzoeksvragen}%
\label{subsec:onderzoeksvragen}

De hoofdonderzoeksvraag luidt: 
“Hoe kan een kostenefficiënte en GDPR-conforme cloudarchitectuur op Microsoft Azure worden ontworpen en als Proof of Concept worden uitgewerkt voor de opslag van Zorglab-simulatiedata, rekening houdend met dataminimalisatie en regelgebaseerde filtering?”

Om deze centrale hoofdvraag te beantwoorden, worden de volgende deelvragen onderzocht:

\subsubsection*{Deelvragen Probleemdomein}
\begin{enumerate}
    \item \textbf{Huidige praktijk \& brondata:} Wat is de huidige operationele workflow en technische opzet van de videorecordings in het Zorglab HOGENT, en welke datavolumes, formaten en technische knelpunten brengt dit met zich mee?
    \item \textbf{Privacy \& GDPR:} Welke wettelijke en ethische privacyverplichtingen gelden er voor de verwerking van videobeelden van studenten en simulanten?
    \item \textbf{Kritische dataminimalisatie:} Hoe kan het principe van dataminimalisatie kritisch worden toegepast om te bepalen welke videodata effectief in aanmerking komt voor verdere verwerking, en welke data omwille van privacy, beperkte meerwaarde of overbodigheid beter niet naar de cloud wordt doorgestuurd?
\end{enumerate}

\subsubsection*{Deelvragen Oplossingsdomein}
\begin{enumerate}
    \setcounter{enumi}{3}
    \item \textbf{Cloudarchitectuur:} Hoe kan een kostenefficiënte en schaalbare cloudarchitectuur op Microsoft Azure worden ontworpen, om de geselecteerde videodata van het Zorglab op te slaan?
    \item \textbf{Automatisering \& PoC:} Hoe kunnen geautomatiseerde opslag- en opschoonregels op Microsoft Azure worden ingesteld om de bewaartermijnen van de videodata GDPR-compliant te beheren binnen een Proof of Concept?
\end{enumerate}

\subsection{Scope van de Proof of Concept (PoC)}%
\label{subsec:scope}

Om de haalbaarheid van het onderzoek te bewaken, richt de praktische uitwerking (PoC) zich op:
\begin{itemize}
    \item \textbf{In scope:} Het opzetten van een Azure Resource Group, Blob Storage Accounts met verschillende storage tiers, het toepassen van metadata-tagging en het configureren van geautomatiseerde Lifecycle Management rules voor transitie en verwijdering van test-videobestanden.
    \item \textbf{Out of scope:} Het ontwikkelen van complexe AI-modellen voor automatische video-analyse of de volledige fysieke herinrichting van het netwerk in het Zorglab.
\end{itemize}

%---------- Stand van zaken ---------------------------------------------------

\section{Literatuurstudie}%
\label{sec:literatuurstudie}

\subsection{Probleemdomein}%
\label{subsec:lit_probleemdomein}

\subsubsection{Technische opzet en datastromen in de simulatiepraktijk}
In moderne medische simulatie-omgevingen -- en bij uitbreiding ook binnen het Zorglab van HOGENT -- wordt gebruikgemaakt van geavanceerde audio- en videosystemen om praktijkoefeningen vast te leggen voor debriefing en evaluatie \parencite{CHSU2023}. De technische opzet combineert vaak softwareoplossingen zoals Noldus VISO met IP-camera's (zoals Axis) en samenwerkingsplatformen zoals Microsoft Teams.

Het voornaamste technische knelpunt bij deze opzet is de generatie van omvangrijke videodata. Hoge resoluties en multi-camerahoeken leiden tot een exponentiële groei van datavolumes op lokale opslagmedia. Zonder geautomatiseerd beheer of tijdig opschonen ontstaan er al snel knelpunten rond opslagcapaciteit, bandbreedte en beheerskosten.

\subsubsection{Juridisch en ethisch kader rond videodata van studenten}
Videobeelden waarop studenten of simulanten herkenbaar in beeld gebracht worden, vallen onder de definitie van persoonsgegevens conform de Algemene Verordening Gegevensbescherming (AVG/GDPR) \parencite{Sensors2025}. Hierdoor geldt een strikt juridisch en ethisch kader voor de verwerking en opslag van deze data.

Belangrijke verplichtingen binnen dit kader omvatten:
\begin{itemize}
    \item \textbf{Verwerkingsgrondslag en transparantie:} Het expliciet informeren van betrokkenen en het vastleggen van een correcte juridische basis voor verwerking \parencite{Chassang2017}.
    \item \textbf{Datasecuriteit en opslaglocatie:} Binnen de Europese Unie moet gegarandeerd worden dat persoonsgegevens worden opgeslagen binnen conforme datacenters met adequate toegangsbeveiliging \parencite{Sensors2025}.
    \item \textbf{Onderscheid onderwijs- versus onderzoeksdata:} Videobeelden die puur dienen voor onderwijsevaluaties (debriefing) kennen een korte bewaartermijn. Beelden die worden ingezet voor wetenschappelijk onderzoek vereisen specifieke toestemming en mogen onder strengere voorwaarden langer bewaard worden \parencite{CHSU2023}.
\end{itemize}

\subsubsection{Toepassing van dataminimalisatie en selectiecriteria}
Het AVG-principe van dataminimalisatie eist dat niet meer gegevens worden verwerkt en bewaard dan strikt noodzakelijk is voor het vooropgestelde doel \parencite{Chassang2017}. In een videorecording-context betekent dit dat niet alle gemaakte opnames onbeperkt naar een centrale opslag of de cloud mogen doorstromen.

Door het instellen van kritische filtercriteria – gebaseerd op metadata zoals het type sessie, de geldigheid van toestemmingsformulieren of de relevantie voor onderzoek – kan aan de bron bepaald worden welke data in aanmerking komt voor opslag en welke data omwille van privacy, beperkte meerwaarde of overbodigheid beter niet naar de cloud wordt doorgestuurd \parencite{CHSU2023}.

\subsection{Oplossingsdomein}%
\label{subsec:lit_oplossingsdomein}

\subsubsection{Cloudarchitectuur en kostenefficiënte opslag op Microsoft Azure}
Om de uitdagingen rond datagroei en veilige opslag op te lossen, biedt een cloudarchitectuur gebaseerd op Microsoft Azure uitkomst \parencite{AzureWellArchitected2026}. Binnen Azure Blob Storage kunnen gegevens kostenefficiënt georganiseerd worden door gebruik te maken van verschillende opslaglagen (\emph{storage tiers}) \parencite{AzureLifecycle2026}:
\begin{itemize}
    \item \textbf{Hot Tier:} Geoptimaliseerd voor data die frequent geraadpleegd wordt \parencite{AzureLifecycle2026}.
    \item \textbf{Cool Tier:} Geschikt voor data die minder vaak geraadpleegd wordt, maar wel direct beschikbaar moet zijn \parencite{AzureLifecycle2026}.
    \item \textbf{Cold / Archive Tier:} Voor langdurige archivering van onderzoeksdata tegen zeer lage opslagkosten \parencite{AzureWellArchitected2026}.
\end{itemize}

\subsubsection{Geautomatiseerd Lifecycle Management en regelgebaseerde opschoning}
Het handmatig beheren en verwijderen van videobestanden is foutgevoelig en tijdsintensief. Microsoft Azure biedt hiervoor \emph{Lifecycle Management rules} \parencite{AzureLifecycle2026}. Dit zijn geautomatiseerde, regelgebaseerde beleidsregels die op basis van paden, containerbronnen of specifieke metadata-tags acties uitvoeren.

Via deze regels worden bestanden na het verstrijken van de vastgelegde bewaartermijn automatisch getransfereerd naar een goedkopere opslaglaag of definitief verwijderd \parencite{AzureLifecycle2026}. Dit garandeert dat bewaartermijnen GDPR-compliant worden afgedwongen zonder menselijke tussenkomst \parencite{AzureWellArchitected2026}.

% Voor literatuurverwijzingen zijn er twee belangrijke commando's:
% \autocite{KEY} => (Auteur, jaartal) Gebruik dit als de naam van de auteur
%   geen onderdeel is van de zin.
% \textcite{KEY} => Auteur (jaartal)  Gebruik dit als de auteursnaam wel een
%   functie heeft in de zin (bv. ``Uit onderzoek door Doll & Hill (1954) bleek
%   ...'')

%---------- Methodologie ------------------------------------------------------
\section{Methodologie}%
\label{sec:methodologie}

Om de hoofdonderzoeksvraag te beantwoorden en tot een kostenefficiënte, GDPR-conforme cloudarchitectuur te komen, wordt een praktijkgericht onderzoek uitgevoerd. De methodologie volgt vier iteratieve fasen:

\subsection{Fase 1: Analyse van het privacykader}
In de eerste fase wordt een grondige analyse uitgevoerd van de wettelijke en ethische randvoorwaarden. Hierbij staan de AVG/GDPR, de vereisten rond EU-datalocatie en de specifieke richtlijnen van de Data Protection Officer (DPO) centraal om de eisen rond bewaartermijnen, toegangsbeheer en datasecuriteit vast te leggen.

\subsection{Fase 2: In kaart brengen van de huidige datastromen}
Vervolgens worden de bestaande operationele workflows binnen het Zorglab geanalyseerd. De datastromen afkomstig van Noldus VISO, Axis-camera's en Microsoft Teams worden gedocumenteerd om inzicht te krijgen in de gegenereerde datavolumes, bestandsformaten en technische knelpunten rond de huidige opslag.

\subsection{Fase 3: Bepalen van regelgebaseerde filterregels}
In deze fase worden de criteria voor dataminimalisatie en opslagbeheer gedefinieerd. Op basis van metadata wordt vastgelegd welke data naar de cloud doorstroomt en hoe deze efficiënt wordt verdeeld over de verschillende Azure-opslaglagen.

\subsection{Fase 4: Ontwikkeling van de Proof of Concept}
Tot slot wordt het ontworpen model in de praktijk gebracht en getest binnen Microsoft Azure. In deze fase worden de Azure Resource Groups, Blob Storage Accounts, metadata-tagging en geautomatiseerde geconfigureerd en gevalideerd met test-videobestanden.


%---------- Verwachte resultaten ----------------------------------------------
\section{Verwacht resultaat, conclusie}
\label{sec:verwachte_resultaten}

Het hoofddoel van deze bachelorproef is het opleveren van een theoretisch en praktisch onderbouwd kader voor de geautomatiseerde, veilige en kostenefficiënte opslag van Zorglab-simulatiedata op Microsoft Azure. Dit resulteert in een analyse van de wettelijke randvoorwaarden gekoppeld aan de huidige datastromen, een regelgebaseerd filtermodel op basis van metadata, en een geteste Proof of Concept op Azure met geconfigureerde.

Hiermee verkrijgt het Zorglab een schaalbaar en toekomstbestendig fundament. De automatische afdwinging van bewaartermijnen en de tijdige migratie naar goedkopere opslaglagen waarborgen de strikte naleving van de privacywetgeving en reduceren de opslagkosten tot het daadwerkelijke gebruik.


\printbibliography[heading=bibintoc]

\end{document}